% Options for packages loaded elsewhere
\PassOptionsToPackage{unicode}{hyperref}
\PassOptionsToPackage{hyphens}{url}
\documentclass[
]{article}
\usepackage{xcolor}
\usepackage[margin=1in]{geometry}
\usepackage{amsmath,amssymb}
\setcounter{secnumdepth}{-\maxdimen} % remove section numbering
\usepackage{iftex}
\ifPDFTeX
  \usepackage[T1]{fontenc}
  \usepackage[utf8]{inputenc}
  \usepackage{textcomp} % provide euro and other symbols
\else % if luatex or xetex
  \usepackage{unicode-math} % this also loads fontspec
  \defaultfontfeatures{Scale=MatchLowercase}
  \defaultfontfeatures[\rmfamily]{Ligatures=TeX,Scale=1}
\fi
\usepackage{lmodern}
\ifPDFTeX\else
  % xetex/luatex font selection
\fi
% Use upquote if available, for straight quotes in verbatim environments
\IfFileExists{upquote.sty}{\usepackage{upquote}}{}
\IfFileExists{microtype.sty}{% use microtype if available
  \usepackage[]{microtype}
  \UseMicrotypeSet[protrusion]{basicmath} % disable protrusion for tt fonts
}{}
\makeatletter
\@ifundefined{KOMAClassName}{% if non-KOMA class
  \IfFileExists{parskip.sty}{%
    \usepackage{parskip}
  }{% else
    \setlength{\parindent}{0pt}
    \setlength{\parskip}{6pt plus 2pt minus 1pt}}
}{% if KOMA class
  \KOMAoptions{parskip=half}}
\makeatother
\usepackage{color}
\usepackage{fancyvrb}
\newcommand{\VerbBar}{|}
\newcommand{\VERB}{\Verb[commandchars=\\\{\}]}
\DefineVerbatimEnvironment{Highlighting}{Verbatim}{commandchars=\\\{\}}
% Add ',fontsize=\small' for more characters per line
\usepackage{framed}
\definecolor{shadecolor}{RGB}{248,248,248}
\newenvironment{Shaded}{\begin{snugshade}}{\end{snugshade}}
\newcommand{\AlertTok}[1]{\textcolor[rgb]{0.94,0.16,0.16}{#1}}
\newcommand{\AnnotationTok}[1]{\textcolor[rgb]{0.56,0.35,0.01}{\textbf{\textit{#1}}}}
\newcommand{\AttributeTok}[1]{\textcolor[rgb]{0.13,0.29,0.53}{#1}}
\newcommand{\BaseNTok}[1]{\textcolor[rgb]{0.00,0.00,0.81}{#1}}
\newcommand{\BuiltInTok}[1]{#1}
\newcommand{\CharTok}[1]{\textcolor[rgb]{0.31,0.60,0.02}{#1}}
\newcommand{\CommentTok}[1]{\textcolor[rgb]{0.56,0.35,0.01}{\textit{#1}}}
\newcommand{\CommentVarTok}[1]{\textcolor[rgb]{0.56,0.35,0.01}{\textbf{\textit{#1}}}}
\newcommand{\ConstantTok}[1]{\textcolor[rgb]{0.56,0.35,0.01}{#1}}
\newcommand{\ControlFlowTok}[1]{\textcolor[rgb]{0.13,0.29,0.53}{\textbf{#1}}}
\newcommand{\DataTypeTok}[1]{\textcolor[rgb]{0.13,0.29,0.53}{#1}}
\newcommand{\DecValTok}[1]{\textcolor[rgb]{0.00,0.00,0.81}{#1}}
\newcommand{\DocumentationTok}[1]{\textcolor[rgb]{0.56,0.35,0.01}{\textbf{\textit{#1}}}}
\newcommand{\ErrorTok}[1]{\textcolor[rgb]{0.64,0.00,0.00}{\textbf{#1}}}
\newcommand{\ExtensionTok}[1]{#1}
\newcommand{\FloatTok}[1]{\textcolor[rgb]{0.00,0.00,0.81}{#1}}
\newcommand{\FunctionTok}[1]{\textcolor[rgb]{0.13,0.29,0.53}{\textbf{#1}}}
\newcommand{\ImportTok}[1]{#1}
\newcommand{\InformationTok}[1]{\textcolor[rgb]{0.56,0.35,0.01}{\textbf{\textit{#1}}}}
\newcommand{\KeywordTok}[1]{\textcolor[rgb]{0.13,0.29,0.53}{\textbf{#1}}}
\newcommand{\NormalTok}[1]{#1}
\newcommand{\OperatorTok}[1]{\textcolor[rgb]{0.81,0.36,0.00}{\textbf{#1}}}
\newcommand{\OtherTok}[1]{\textcolor[rgb]{0.56,0.35,0.01}{#1}}
\newcommand{\PreprocessorTok}[1]{\textcolor[rgb]{0.56,0.35,0.01}{\textit{#1}}}
\newcommand{\RegionMarkerTok}[1]{#1}
\newcommand{\SpecialCharTok}[1]{\textcolor[rgb]{0.81,0.36,0.00}{\textbf{#1}}}
\newcommand{\SpecialStringTok}[1]{\textcolor[rgb]{0.31,0.60,0.02}{#1}}
\newcommand{\StringTok}[1]{\textcolor[rgb]{0.31,0.60,0.02}{#1}}
\newcommand{\VariableTok}[1]{\textcolor[rgb]{0.00,0.00,0.00}{#1}}
\newcommand{\VerbatimStringTok}[1]{\textcolor[rgb]{0.31,0.60,0.02}{#1}}
\newcommand{\WarningTok}[1]{\textcolor[rgb]{0.56,0.35,0.01}{\textbf{\textit{#1}}}}
\usepackage{longtable,booktabs,array}
\usepackage{calc} % for calculating minipage widths
% Correct order of tables after \paragraph or \subparagraph
\usepackage{etoolbox}
\makeatletter
\patchcmd\longtable{\par}{\if@noskipsec\mbox{}\fi\par}{}{}
\makeatother
% Allow footnotes in longtable head/foot
\IfFileExists{footnotehyper.sty}{\usepackage{footnotehyper}}{\usepackage{footnote}}
\makesavenoteenv{longtable}
\usepackage{graphicx}
\makeatletter
\newsavebox\pandoc@box
\newcommand*\pandocbounded[1]{% scales image to fit in text height/width
  \sbox\pandoc@box{#1}%
  \Gscale@div\@tempa{\textheight}{\dimexpr\ht\pandoc@box+\dp\pandoc@box\relax}%
  \Gscale@div\@tempb{\linewidth}{\wd\pandoc@box}%
  \ifdim\@tempb\p@<\@tempa\p@\let\@tempa\@tempb\fi% select the smaller of both
  \ifdim\@tempa\p@<\p@\scalebox{\@tempa}{\usebox\pandoc@box}%
  \else\usebox{\pandoc@box}%
  \fi%
}
% Set default figure placement to htbp
\def\fps@figure{htbp}
\makeatother
\setlength{\emergencystretch}{3em} % prevent overfull lines
\providecommand{\tightlist}{%
  \setlength{\itemsep}{0pt}\setlength{\parskip}{0pt}}
\usepackage{bookmark}
\IfFileExists{xurl.sty}{\usepackage{xurl}}{} % add URL line breaks if available
\urlstyle{same}
\hypersetup{
  hidelinks,
  pdfcreator={LaTeX via pandoc}}

\author{}
\date{\vspace{-2.5em}}

\begin{document}

\section{Qemu Calypso}\label{qemu-calypso}

Ce document décrit l'émulation complète du SoC \textbf{TI Calypso DBB}
(Digital Baseband) dans QEMU, permettant d'exécuter les firmwares
OsmocomBB (loader, layer1) sans matériel physique (Motorola C123/C118).

Le Calypso est un ARM946E-S avec un ensemble de périphériques
spécifiques à la téléphonie GSM. Aucune émulation publique complète
n'existait avant ce travail.

\begin{center}\rule{0.5\linewidth}{0.5pt}\end{center}

\subsection{Architecture globale}\label{architecture-globale}

\paragraph{Vue d'ensemble de la pile GSM
émulée}\label{vue-densemble-de-la-pile-gsm-uxe9muluxe9e}

\begin{Shaded}
\begin{Highlighting}[]
\NormalTok{dot\_code }\OtherTok{\textless{}{-}} \StringTok{\textquotesingle{}digraph gsm\_stack \{}
\StringTok{  graph [rankdir=TB, bgcolor=transparent, fontname=Helvetica, pad="0.5"]}
\StringTok{  node [shape=box, style="filled,rounded", fontname=Helvetica, fontsize=11]}
\StringTok{  edge [fontname=Helvetica, fontsize=9]}

\StringTok{  subgraph cluster\_host \{}
\StringTok{    label="Host Linux"}
\StringTok{    style=dashed; color="\#666666"; fontcolor="\#666666"}

\StringTok{    wireshark [label="Wireshark}\SpecialCharTok{\textbackslash{}n}\StringTok{(GSMTAP)", fillcolor="\#E8F5E9"]}
\StringTok{    osmocon [label="osmocon}\SpecialCharTok{\textbackslash{}n}\StringTok{(sercomm bridge)", fillcolor="\#FFF9C4"]}
\StringTok{    mobile [label="mobile / cell\_log}\SpecialCharTok{\textbackslash{}n}\StringTok{(layer23)", fillcolor="\#FFF9C4"]}
\StringTok{    nitb [label="osmo{-}nitb}\SpecialCharTok{\textbackslash{}n}\StringTok{(MSC+HLR+BSC)", fillcolor="\#BBDEFB"]}
\StringTok{    bts [label="osmo{-}bts{-}trx}\SpecialCharTok{\textbackslash{}n}\StringTok{(BTS)", fillcolor="\#BBDEFB"]}
\StringTok{    transceiver [label="transceiver}\SpecialCharTok{\textbackslash{}n}\StringTok{(L1→TRX bridge)", fillcolor="\#FFE0B2"]}
\StringTok{  \}}

\StringTok{  subgraph cluster\_qemu \{}
\StringTok{    label="QEMU — TI Calypso SoC"}
\StringTok{    style=filled; fillcolor="\#F3E5F5"; color="\#7B1FA2"; fontcolor="\#7B1FA2"}

\StringTok{    firmware [label="layer1.highram.elf}\SpecialCharTok{\textbackslash{}n}\StringTok{(OsmocomBB L1)", fillcolor="\#CE93D8"]}
\StringTok{    uart [label="UART Modem}\SpecialCharTok{\textbackslash{}n}\StringTok{0xFFFF5800", fillcolor="\#E1BEE7"]}
\StringTok{    trx [label="TRX Bridge}\SpecialCharTok{\textbackslash{}n}\StringTok{(UDP 4729)", fillcolor="\#E1BEE7"]}
\StringTok{    inth [label="INTH}\SpecialCharTok{\textbackslash{}n}\StringTok{0xFFFFFA00", fillcolor="\#E1BEE7"]}
\StringTok{  \}}

\StringTok{  firmware {-}\textgreater{} uart [label="sercomm}\SpecialCharTok{\textbackslash{}n}\StringTok{frames"]}
\StringTok{  uart {-}\textgreater{} osmocon [label="pty}\SpecialCharTok{\textbackslash{}n}\StringTok{(serial)"]}
\StringTok{  osmocon {-}\textgreater{} mobile [label="L1CTL}\SpecialCharTok{\textbackslash{}n}\StringTok{/tmp/osmocom\_l2"]}
\StringTok{  transceiver {-}\textgreater{} trx [label="UDP}\SpecialCharTok{\textbackslash{}n}\StringTok{4729"]}
\StringTok{  transceiver {-}\textgreater{} osmocon [label="L1CTL", style=dashed]}
\StringTok{  bts {-}\textgreater{} transceiver [label="TRX}\SpecialCharTok{\textbackslash{}n}\StringTok{protocol"]}
\StringTok{  bts {-}\textgreater{} nitb [label="Abis/IP"]}
\StringTok{  nitb {-}\textgreater{} wireshark [label="GSMTAP}\SpecialCharTok{\textbackslash{}n}\StringTok{UDP 4729", style=dashed, color="\#4CAF50"]}
\StringTok{  firmware {-}\textgreater{} trx [label="DSP API}\SpecialCharTok{\textbackslash{}n}\StringTok{(shared RAM)"]}
\StringTok{  firmware {-}\textgreater{} inth [label="IRQs"]}
\StringTok{\}\textquotesingle{}}
\FunctionTok{render\_graphviz}\NormalTok{(dot\_code)}
\end{Highlighting}
\end{Shaded}

\begin{center}\includegraphics[width=4.51in]{../../../../tmp/RtmpZyT17y/file2730a5eec7a3a} \end{center}

\paragraph{Architecture interne du SoC Calypso
émulé}\label{architecture-interne-du-soc-calypso-uxe9muluxe9}

\begin{Shaded}
\begin{Highlighting}[]
\NormalTok{dot\_code }\OtherTok{\textless{}{-}} \StringTok{\textquotesingle{}digraph calypso\_soc \{}
\StringTok{  graph [rankdir=LR, bgcolor=transparent, fontname=Helvetica]}
\StringTok{  node [shape=record, style=filled, fontname="Courier", fontsize=10]}
\StringTok{  edge [fontname=Helvetica, fontsize=9]}

\StringTok{  cpu [label="\{ARM946E{-}S|PC, CPSR, R0{-}R15|Modes: SVC/IRQ/FIQ/ABT/UND\}",}
\StringTok{       fillcolor="\#FFCDD2"]}

\StringTok{  inth [label="\{INTH|0xFFFFFA00|32 IRQ lines|Mask, Pending, Priority\}",}
\StringTok{        fillcolor="\#F8BBD0"]}

\StringTok{  subgraph cluster\_periphs \{}
\StringTok{    label="Peripherals"; style=dashed; color="\#999"}

\StringTok{    timer1 [label="\{Timer 1|0xFFFE3800|IRQ 1\}", fillcolor="\#C8E6C9"]}
\StringTok{    timer2 [label="\{Timer 2|0xFFFE3C00|IRQ 2\}", fillcolor="\#C8E6C9"]}
\StringTok{    uart\_m [label="\{UART Modem|0xFFFF5800|IRQ 7|16550{-}like\}", fillcolor="\#B3E5FC"]}
\StringTok{    uart\_i [label="\{UART IrDA|0xFFFF5000|IRQ 18\}", fillcolor="\#B3E5FC"]}
\StringTok{    spi    [label="\{SPI (ABB)|0xFFFE3000|IRQ 13|TWL3025\}", fillcolor="\#DCEDC8"]}
\StringTok{    i2c    [label="\{I2C stub|0xFFFE1800\}", fillcolor="\#F0F4C3"]}
\StringTok{    keypad [label="\{Keypad|0xFFFE4800|Active{-}low|POWER pulse\}", fillcolor="\#FFE0B2"]}
\StringTok{    trx    [label="\{TRX Bridge|TPU+DSP API|UDP 4729\}", fillcolor="\#D1C4E9"]}
\StringTok{  \}}

\StringTok{  memory [label="\{Memory Map|IRAM 0x00800000 (256K)|Alias → 0x00000000|Catch{-}all 0xFFF00000\}",}
\StringTok{          fillcolor="\#FFF9C4"]}

\StringTok{  cpu {-}\textgreater{} inth [label="IRQ/FIQ"]}
\StringTok{  inth {-}\textgreater{} timer1 [label="IRQ 1", dir=back]}
\StringTok{  inth {-}\textgreater{} timer2 [label="IRQ 2", dir=back]}
\StringTok{  inth {-}\textgreater{} uart\_m [label="IRQ 7", dir=back]}
\StringTok{  inth {-}\textgreater{} uart\_i [label="IRQ 18", dir=back]}
\StringTok{  inth {-}\textgreater{} spi [label="IRQ 13", dir=back]}
\StringTok{  cpu {-}\textgreater{} memory [label="bus"]}
\StringTok{\}\textquotesingle{}}
\FunctionTok{render\_graphviz}\NormalTok{(dot\_code)}
\end{Highlighting}
\end{Shaded}

\begin{center}\includegraphics[width=9in]{../../../../tmp/RtmpZyT17y/file2730a7c543b09} \end{center}

\begin{center}\rule{0.5\linewidth}{0.5pt}\end{center}

\subsection{Carte mémoire}\label{carte-muxe9moire}

\begin{Shaded}
\begin{Highlighting}[]
\NormalTok{dot\_code }\OtherTok{\textless{}{-}} \StringTok{\textquotesingle{}digraph memmap \{}
\StringTok{  graph [rankdir=TB, bgcolor=transparent, fontname=Helvetica]}
\StringTok{  node [shape=record, style=filled, fontname="Courier", fontsize=9,}
\StringTok{        fillcolor="\#E3F2FD"]}
\StringTok{  edge [style=invis]}

\StringTok{  title [label="Calypso Memory Map", shape=plaintext,}
\StringTok{         fontsize=14, fontname="Helvetica Bold"]}

\StringTok{  m0 [label="\{0x00000000|IRAM Alias (256K)|Miroir de 0x00800000\}"]}
\StringTok{  m1 [label="\{0x00000300|Stub MMIO (16{-}bit)\}"]}
\StringTok{  m2 [label="\{0x00800000|IRAM (256K)|Code + Data|loader / layer1 @ 0x00820000\}",}
\StringTok{      fillcolor="\#C8E6C9"]}
\StringTok{  m3 [label="\{0xFFFE1800|I2C stub\}"]}
\StringTok{  m4 [label="\{0xFFFE3000|SPI (ABB TWL3025)\}"]}
\StringTok{  m5 [label="\{0xFFFE3800|Timer 1\}"]}
\StringTok{  m6 [label="\{0xFFFE3C00|Timer 2\}"]}
\StringTok{  m7 [label="\{0xFFFE4800|Keypad (active{-}low)\}"]}
\StringTok{  m8 [label="\{0xFFFF5000|UART IrDA\}"]}
\StringTok{  m9 [label="\{0xFFFF5800|UART Modem (osmocon)\}",}
\StringTok{      fillcolor="\#FFECB3"]}
\StringTok{  m10 [label="\{0xFFFFF000|DPLL stub\}"]}
\StringTok{  m11 [label="\{0xFFFFFA00|INTH (Interrupt Handler)\}",}
\StringTok{       fillcolor="\#F8BBD0"]}
\StringTok{  m12 [label="\{0xFFFFFB00|CLKM stub\}"]}
\StringTok{  m13 [label="\{0xFFF00000|Catch{-}all (1MB, lowest prio)\}",}
\StringTok{       fillcolor="\#FFCDD2"]}

\StringTok{  title {-}\textgreater{} m0 {-}\textgreater{} m1 {-}\textgreater{} m2 {-}\textgreater{} m3 {-}\textgreater{} m4 {-}\textgreater{} m5 {-}\textgreater{} m6 {-}\textgreater{} m7 {-}\textgreater{} m8 {-}\textgreater{} m9 {-}\textgreater{} m10 {-}\textgreater{} m11 {-}\textgreater{} m12 {-}\textgreater{} m13}
\StringTok{\}\textquotesingle{}}
\FunctionTok{render\_graphviz}\NormalTok{(dot\_code)}
\end{Highlighting}
\end{Shaded}

\begin{center}\includegraphics[width=1.19in]{../../../../tmp/RtmpZyT17y/file2730a2bae9bdd} \end{center}

\begin{center}\rule{0.5\linewidth}{0.5pt}\end{center}

\subsection{Périphériques
implémentés}\label{puxe9riphuxe9riques-impluxe9mentuxe9s}

\paragraph{INTH --- Interrupt Controller
(0xFFFFFA00)}\label{inth-interrupt-controller-0xfffffa00}

Le contrôleur d'interruptions Calypso gère 32 lignes d'IRQ avec masquage
et priorités.

\begin{Shaded}
\begin{Highlighting}[]
\NormalTok{dot\_code }\OtherTok{\textless{}{-}} \StringTok{\textquotesingle{}digraph inth \{}
\StringTok{  graph [rankdir=LR, bgcolor=transparent, fontname=Helvetica]}
\StringTok{  node [shape=box, style="filled,rounded", fontname=Helvetica, fontsize=10]}
\StringTok{  edge [fontname=Helvetica, fontsize=9]}

\StringTok{  subgraph cluster\_sources \{}
\StringTok{    label="IRQ Sources"; style=dashed}
\StringTok{    t1 [label="Timer1}\SpecialCharTok{\textbackslash{}n}\StringTok{IRQ 1", fillcolor="\#C8E6C9"]}
\StringTok{    t2 [label="Timer2}\SpecialCharTok{\textbackslash{}n}\StringTok{IRQ 2", fillcolor="\#C8E6C9"]}
\StringTok{    uart [label="UART Modem}\SpecialCharTok{\textbackslash{}n}\StringTok{IRQ 7", fillcolor="\#B3E5FC"]}
\StringTok{    spi [label="SPI/ABB}\SpecialCharTok{\textbackslash{}n}\StringTok{IRQ 13", fillcolor="\#DCEDC8"]}
\StringTok{    irda [label="UART IrDA}\SpecialCharTok{\textbackslash{}n}\StringTok{IRQ 18", fillcolor="\#B3E5FC"]}
\StringTok{    tdma [label="TDMA Frame}\SpecialCharTok{\textbackslash{}n}\StringTok{IRQ (TRX)", fillcolor="\#D1C4E9"]}
\StringTok{  \}}

\StringTok{  subgraph cluster\_inth \{}
\StringTok{    label="INTH Registers"; style=filled; fillcolor="\#FCE4EC"}
\StringTok{    itr [label="ITR}\SpecialCharTok{\textbackslash{}n}\StringTok{(Pending)", shape=record]}
\StringTok{    mir [label="MIR}\SpecialCharTok{\textbackslash{}n}\StringTok{(Mask)", shape=record]}
\StringTok{    ilr [label="ILR0{-}31}\SpecialCharTok{\textbackslash{}n}\StringTok{(Priority+FIQ)", shape=record]}
\StringTok{    irq\_out [label="IRQ out", shape=diamond, fillcolor="\#EF9A9A"]}
\StringTok{    fiq\_out [label="FIQ out", shape=diamond, fillcolor="\#EF9A9A"]}
\StringTok{  \}}

\StringTok{  cpu [label="ARM946E{-}S}\SpecialCharTok{\textbackslash{}n}\StringTok{CPU", fillcolor="\#FFCDD2"]}

\StringTok{  t1 {-}\textgreater{} itr}
\StringTok{  t2 {-}\textgreater{} itr}
\StringTok{  uart {-}\textgreater{} itr}
\StringTok{  spi {-}\textgreater{} itr}
\StringTok{  irda {-}\textgreater{} itr}
\StringTok{  tdma {-}\textgreater{} itr}
\StringTok{  itr {-}\textgreater{} mir [label="\& \textasciitilde{}mask"]}
\StringTok{  mir {-}\textgreater{} ilr [label="priority}\SpecialCharTok{\textbackslash{}n}\StringTok{select"]}
\StringTok{  ilr {-}\textgreater{} irq\_out [label="IRQ"]}
\StringTok{  ilr {-}\textgreater{} fiq\_out [label="FIQ"]}
\StringTok{  irq\_out {-}\textgreater{} cpu}
\StringTok{  fiq\_out {-}\textgreater{} cpu}
\StringTok{\}\textquotesingle{}}
\FunctionTok{render\_graphviz}\NormalTok{(dot\_code)}
\end{Highlighting}
\end{Shaded}

\begin{center}\includegraphics[width=5.63in]{../../../../tmp/RtmpZyT17y/file2730a5f0a073f} \end{center}

\textbf{Registres clés :}

\begin{longtable}[]{@{}lll@{}}
\toprule\noalign{}
Offset & Registre & Description \\
\midrule\noalign{}
\endhead
\bottomrule\noalign{}
\endlastfoot
0x00 & ITR & Interrupts pending (raw) \\
0x04 & MIR & Mask --- bit=1 masque l'IRQ \\
0x08 & SIR\_IRQ\_CODE & Numéro de l'IRQ la plus prioritaire \\
0x0C & SIR\_FIQ\_CODE & Numéro de la FIQ la plus prioritaire \\
0x10 & CTRL & Contrôle global (new IRQ agreement) \\
0x20+ & ILR{[}0..31{]} & Priority + FIQ/IRQ select par ligne \\
\end{longtable}

\paragraph{UART --- 16550-like (Modem +
IrDA)}\label{uart-16550-like-modem-irda}

\begin{Shaded}
\begin{Highlighting}[]
\NormalTok{dot\_code }\OtherTok{\textless{}{-}} \StringTok{\textquotesingle{}digraph uart\_flow \{}
\StringTok{  graph [rankdir=LR, bgcolor=transparent, fontname=Helvetica]}
\StringTok{  node [shape=box, style="filled,rounded", fontname=Helvetica, fontsize=10]}
\StringTok{  edge [fontname=Helvetica, fontsize=9]}

\StringTok{  subgraph cluster\_uart \{}
\StringTok{    label="UART Modem — 0xFFFF5800"; style=filled; fillcolor="\#E3F2FD"}

\StringTok{    regs [label="Registers}\SpecialCharTok{\textbackslash{}n}\StringTok{THR/RHR/IER/IIR}\SpecialCharTok{\textbackslash{}n}\StringTok{LCR/MCR/LSR/MSR}\SpecialCharTok{\textbackslash{}n}\StringTok{DLL/DLH", shape=record,}
\StringTok{          fillcolor="\#BBDEFB"]}
\StringTok{    txfifo [label="TX FIFO}\SpecialCharTok{\textbackslash{}n}\StringTok{(64 bytes)", fillcolor="\#B3E5FC"]}
\StringTok{    rxfifo [label="RX FIFO}\SpecialCharTok{\textbackslash{}n}\StringTok{(64 bytes)", fillcolor="\#B3E5FC"]}
\StringTok{    irq [label="IRQ}\SpecialCharTok{\textbackslash{}n}\StringTok{logic", shape=diamond, fillcolor="\#81D4FA"]}
\StringTok{  \}}

\StringTok{  firmware [label="OsmocomBB}\SpecialCharTok{\textbackslash{}n}\StringTok{firmware", fillcolor="\#CE93D8"]}
\StringTok{  chardev [label="QEMU chardev}\SpecialCharTok{\textbackslash{}n}\StringTok{(pty)", fillcolor="\#FFF9C4"]}
\StringTok{  osmocon [label="osmocon}\SpecialCharTok{\textbackslash{}n}\StringTok{(host)", fillcolor="\#FFE0B2"]}
\StringTok{  inth [label="INTH}\SpecialCharTok{\textbackslash{}n}\StringTok{IRQ 7", fillcolor="\#F8BBD0"]}

\StringTok{  firmware {-}\textgreater{} regs [label="MMIO}\SpecialCharTok{\textbackslash{}n}\StringTok{write THR"]}
\StringTok{  regs {-}\textgreater{} txfifo [label="TX"]}
\StringTok{  txfifo {-}\textgreater{} chardev [label="chr\_fe\_write"]}
\StringTok{  chardev {-}\textgreater{} osmocon [label="pty bytes"]}
\StringTok{  osmocon {-}\textgreater{} chardev [label="pty bytes", style=dashed]}
\StringTok{  chardev {-}\textgreater{} rxfifo [label="chr\_fe\_handler"]}
\StringTok{  rxfifo {-}\textgreater{} regs [label="RX"]}
\StringTok{  regs {-}\textgreater{} firmware [label="MMIO}\SpecialCharTok{\textbackslash{}n}\StringTok{read RHR", style=dashed]}
\StringTok{  irq {-}\textgreater{} inth [label="pulse"]}
\StringTok{  regs {-}\textgreater{} irq [label="IER}\SpecialCharTok{\textbackslash{}n}\StringTok{conditions"]}
\StringTok{\}\textquotesingle{}}
\FunctionTok{render\_graphviz}\NormalTok{(dot\_code)}
\end{Highlighting}
\end{Shaded}

\begin{center}\includegraphics[width=6.35in]{../../../../tmp/RtmpZyT17y/file2730a14c8c475} \end{center}

\textbf{Points clés de l'implémentation :}

\begin{itemize}
\tightlist
\item
  Le chardev est attaché via \texttt{qdev\_prop\_set\_chr()}
  \textbf{avant} \texttt{realize()}
\item
  Les handlers RX sont configurés dans \texttt{calypso\_uart\_realize()}
\item
  Protocole sercomm encapsulé dans les trames UART
\item
  Le loader utilise le protocole HDLC pour le handshake osmocon
\item
  Le layer1 utilise sercomm pour L1CTL
\end{itemize}

\paragraph{Timers (Timer1 @ 0xFFFE3800, Timer2 @
0xFFFE3C00)}\label{timers-timer1-0xfffe3800-timer2-0xfffe3c00}

\begin{Shaded}
\begin{Highlighting}[]
\NormalTok{dot\_code }\OtherTok{\textless{}{-}} \StringTok{\textquotesingle{}digraph timer \{}
\StringTok{  graph [rankdir=TB, bgcolor=transparent, fontname=Helvetica]}
\StringTok{  node [shape=record, style=filled, fontname="Courier", fontsize=10]}
\StringTok{  edge [fontname=Helvetica, fontsize=9]}

\StringTok{  regs [label="\{Timer Registers|CNTL (mode, start/stop)|LOAD (reload value)|READ (current count)\}",}
\StringTok{        fillcolor="\#C8E6C9"]}
\StringTok{  qemu\_timer [label="\{QEMUTimer|QEMU\_CLOCK\_VIRTUAL|callback on expiry\}",}
\StringTok{               fillcolor="\#FFF9C4"]}
\StringTok{  irq [label="IRQ → INTH", shape=diamond, fillcolor="\#F8BBD0"]}

\StringTok{  regs {-}\textgreater{} qemu\_timer [label="timer\_mod()}\SpecialCharTok{\textbackslash{}n}\StringTok{on CNTL write"]}
\StringTok{  qemu\_timer {-}\textgreater{} irq [label="fire on}\SpecialCharTok{\textbackslash{}n}\StringTok{expiry"]}
\StringTok{  qemu\_timer {-}\textgreater{} regs [label="auto{-}reload}\SpecialCharTok{\textbackslash{}n}\StringTok{if AR bit set", style=dashed]}
\StringTok{\}\textquotesingle{}}
\FunctionTok{render\_graphviz}\NormalTok{(dot\_code)}
\end{Highlighting}
\end{Shaded}

\begin{center}\includegraphics[width=1.13in]{../../../../tmp/RtmpZyT17y/file2730a3ff986ab} \end{center}

\paragraph{SPI / ABB TWL3025}\label{spi-abb-twl3025}

Le SPI communique avec l'ABB (Analog Baseband) TWL3025 qui gère la
radio, le VCXO, l'ADC batterie, etc.

\begin{Shaded}
\begin{Highlighting}[]
\NormalTok{dot\_code }\OtherTok{\textless{}{-}} \StringTok{\textquotesingle{}digraph spi\_abb \{}
\StringTok{  graph [rankdir=LR, bgcolor=transparent, fontname=Helvetica]}
\StringTok{  node [shape=box, style="filled,rounded", fontname=Helvetica, fontsize=10]}

\StringTok{  spi [label="SPI Controller}\SpecialCharTok{\textbackslash{}n}\StringTok{0xFFFE3000}\SpecialCharTok{\textbackslash{}n}\StringTok{IRQ 13", fillcolor="\#DCEDC8"]}
\StringTok{  abb [label="ABB TWL3025}\SpecialCharTok{\textbackslash{}n}\StringTok{(émulé en registres)}\SpecialCharTok{\textbackslash{}n}\StringTok{VRPCDEV=0x01}\SpecialCharTok{\textbackslash{}n}\StringTok{VRPCSTEN=0x13",}
\StringTok{       fillcolor="\#FFF9C4"]}
\StringTok{  firmware [label="Firmware}\SpecialCharTok{\textbackslash{}n}\StringTok{abb\_reg\_write()}\SpecialCharTok{\textbackslash{}n}\StringTok{abb\_reg\_read()", fillcolor="\#CE93D8"]}

\StringTok{  firmware {-}\textgreater{} spi [label="MMIO"]}
\StringTok{  spi {-}\textgreater{} abb [label="page/addr/data}\SpecialCharTok{\textbackslash{}n}\StringTok{10{-}bit words"]}
\StringTok{  abb {-}\textgreater{} spi [label="read data", style=dashed]}
\StringTok{\}\textquotesingle{}}
\FunctionTok{render\_graphviz}\NormalTok{(dot\_code)}
\end{Highlighting}
\end{Shaded}

\begin{center}\includegraphics[width=3.66in]{../../../../tmp/RtmpZyT17y/file2730a2f9a084c} \end{center}

\paragraph{Keypad (0xFFFE4800)}\label{keypad-0xfffe4800}

\begin{Shaded}
\begin{Highlighting}[]
\NormalTok{dot\_code }\OtherTok{\textless{}{-}} \StringTok{\textquotesingle{}digraph keypad \{}
\StringTok{  graph [rankdir=TB, bgcolor=transparent, fontname=Helvetica]}
\StringTok{  node [shape=box, style="filled,rounded", fontname=Helvetica, fontsize=10]}
\StringTok{  edge [fontname=Helvetica, fontsize=9]}

\StringTok{  monitor [label="QEMU Monitor}\SpecialCharTok{\textbackslash{}n}\StringTok{write 0xFFFE4800 0xDEAD}\SpecialCharTok{\textbackslash{}n}\StringTok{write 0xFFFE4800 0xBEEF",}
\StringTok{           fillcolor="\#FFF9C4"]}
\StringTok{  kp [label="CalypsoKeypad}\SpecialCharTok{\textbackslash{}n}\StringTok{keys: uint16 (active{-}low)}\SpecialCharTok{\textbackslash{}n}\StringTok{0xFF = all released}\SpecialCharTok{\textbackslash{}n}\StringTok{bit clear = pressed",}
\StringTok{      fillcolor="\#FFE0B2"]}
\StringTok{  timer\_r [label="QEMUTimer}\SpecialCharTok{\textbackslash{}n}\StringTok{release (500ms)", fillcolor="\#C8E6C9"]}
\StringTok{  timer\_2p [label="QEMUTimer}\SpecialCharTok{\textbackslash{}n}\StringTok{2nd press (+700ms)", fillcolor="\#C8E6C9"]}
\StringTok{  timer\_2r [label="QEMUTimer}\SpecialCharTok{\textbackslash{}n}\StringTok{2nd release (+1200ms)", fillcolor="\#C8E6C9"]}
\StringTok{  firmware [label="Firmware}\SpecialCharTok{\textbackslash{}n}\StringTok{keypad scan}\SpecialCharTok{\textbackslash{}n}\StringTok{read 0xFFFE4800", fillcolor="\#CE93D8"]}

\StringTok{  monitor {-}\textgreater{} kp [label="magic write"]}
\StringTok{  kp {-}\textgreater{} timer\_r [label="0xDEAD}\SpecialCharTok{\textbackslash{}n}\StringTok{single press"]}
\StringTok{  kp {-}\textgreater{} timer\_2p [label="0xBEEF}\SpecialCharTok{\textbackslash{}n}\StringTok{double press"]}
\StringTok{  timer\_r {-}\textgreater{} kp [label="set bit}\SpecialCharTok{\textbackslash{}n}\StringTok{(release)"]}
\StringTok{  timer\_2p {-}\textgreater{} kp [label="clear bit}\SpecialCharTok{\textbackslash{}n}\StringTok{(2nd press)"]}
\StringTok{  timer\_2p {-}\textgreater{} timer\_2r [label="schedule"]}
\StringTok{  timer\_2r {-}\textgreater{} kp [label="set bit}\SpecialCharTok{\textbackslash{}n}\StringTok{(2nd release)"]}
\StringTok{  kp {-}\textgreater{} firmware [label="MMIO read", style=dashed]}
\StringTok{\}\textquotesingle{}}
\FunctionTok{render\_graphviz}\NormalTok{(dot\_code)}
\end{Highlighting}
\end{Shaded}

\begin{center}\includegraphics[width=3.17in]{../../../../tmp/RtmpZyT17y/file2730a71332e23} \end{center}

\textbf{Séquence double appui (0xBEEF) :}

\begin{verbatim}
t=0ms         t=500ms       t=700ms        t=1200ms
│─ press 1 ──│── gap ──────│── press 2 ───│
keys=0xFE     keys=0xFF     keys=0xFE      keys=0xFF
\end{verbatim}

\paragraph{TRX Bridge --- DSP API \& TPU}\label{trx-bridge-dsp-api-tpu}

\begin{Shaded}
\begin{Highlighting}[]
\NormalTok{dot\_code }\OtherTok{\textless{}{-}} \StringTok{\textquotesingle{}digraph trx\_bridge \{}
\StringTok{  graph [rankdir=TB, bgcolor=transparent, fontname=Helvetica]}
\StringTok{  node [shape=box, style="filled,rounded", fontname=Helvetica, fontsize=10]}
\StringTok{  edge [fontname=Helvetica, fontsize=9]}

\StringTok{  subgraph cluster\_qemu \{}
\StringTok{    label="QEMU Calypso"; style=filled; fillcolor="\#F3E5F5"}

\StringTok{    tpu [label="TPU Emulation}\SpecialCharTok{\textbackslash{}n}\StringTok{RAM: 0xFFFE8000}\SpecialCharTok{\textbackslash{}n}\StringTok{Reg: 0xFFFF1000", fillcolor="\#D1C4E9"]}
\StringTok{    dsp [label="DSP API}\SpecialCharTok{\textbackslash{}n}\StringTok{Shared RAM}\SpecialCharTok{\textbackslash{}n}\StringTok{0xFFD00000", fillcolor="\#D1C4E9"]}
\StringTok{    trx [label="calypso\_trx}\SpecialCharTok{\textbackslash{}n}\StringTok{UDP socket}\SpecialCharTok{\textbackslash{}n}\StringTok{port 4729", fillcolor="\#CE93D8"]}
\StringTok{    irq [label="TDMA Frame IRQ", shape=diamond, fillcolor="\#F8BBD0"]}
\StringTok{  \}}

\StringTok{  subgraph cluster\_host \{}
\StringTok{    label="Host"; style=dashed; color="\#666"}
\StringTok{    transceiver [label="transceiver}\SpecialCharTok{\textbackslash{}n}\StringTok{(layer23)", fillcolor="\#FFE0B2"]}
\StringTok{    bts [label="osmo{-}bts{-}trx", fillcolor="\#BBDEFB"]}
\StringTok{  \}}

\StringTok{  tpu {-}\textgreater{} trx [label="burst data"]}
\StringTok{  dsp {-}\textgreater{} trx [label="DSP params"]}
\StringTok{  trx {-}\textgreater{} transceiver [label="TRXD}\SpecialCharTok{\textbackslash{}n}\StringTok{UDP 5700+"]}
\StringTok{  transceiver {-}\textgreater{} bts [label="TRX}\SpecialCharTok{\textbackslash{}n}\StringTok{protocol"]}
\StringTok{  trx {-}\textgreater{} irq [label="TDMA tick"]}
\StringTok{  irq {-}\textgreater{} tpu [label="schedule}\SpecialCharTok{\textbackslash{}n}\StringTok{next frame", style=dashed]}
\StringTok{\}\textquotesingle{}}
\FunctionTok{render\_graphviz}\NormalTok{(dot\_code)}
\end{Highlighting}
\end{Shaded}

\begin{center}\includegraphics[width=2.4in]{../../../../tmp/RtmpZyT17y/file2730a5f9d3dee} \end{center}

\begin{center}\rule{0.5\linewidth}{0.5pt}\end{center}

\subsection{Boot sequence}\label{boot-sequence}

\paragraph{Flux de démarrage}\label{flux-de-duxe9marrage}

\begin{Shaded}
\begin{Highlighting}[]
\NormalTok{dot\_code }\OtherTok{\textless{}{-}} \StringTok{\textquotesingle{}digraph boot \{}
\StringTok{  graph [rankdir=TB, bgcolor=transparent, fontname=Helvetica, nodesep=0.3]}
\StringTok{  node [shape=box, style="filled,rounded", fontname=Helvetica, fontsize=10]}
\StringTok{  edge [fontname=Helvetica, fontsize=9]}

\StringTok{  start [label="QEMU start}\SpecialCharTok{\textbackslash{}n}\StringTok{{-}kernel layer1.highram.elf}\SpecialCharTok{\textbackslash{}n}\StringTok{{-}serial pty", fillcolor="\#FFCDD2"]}

\StringTok{  load [label="QEMU loads ELF}\SpecialCharTok{\textbackslash{}n}\StringTok{at 0x00820000}\SpecialCharTok{\textbackslash{}n}\StringTok{(entry point = \_start)", fillcolor="\#FFE0B2"]}

\StringTok{  bss [label="\_start:}\SpecialCharTok{\textbackslash{}n}\StringTok{Clear BSS}\SpecialCharTok{\textbackslash{}n}\StringTok{(loop\_bss)", fillcolor="\#FFF9C4"]}

\StringTok{  stacks [label="Init stacks}\SpecialCharTok{\textbackslash{}n}\StringTok{FIQ sp, IRQ sp}\SpecialCharTok{\textbackslash{}n}\StringTok{SVC sp, ABT sp", fillcolor="\#FFF9C4"]}

\StringTok{  ctors [label="do\_global\_ctors()}\SpecialCharTok{\textbackslash{}n}\StringTok{C++ init", fillcolor="\#FFF9C4"]}

\StringTok{  main [label="\_jump\_main}\SpecialCharTok{\textbackslash{}n}\StringTok{→ main()", fillcolor="\#C8E6C9"]}

\StringTok{  hw\_init [label="Hardware init:}\SpecialCharTok{\textbackslash{}n}\StringTok{UART, SPI/ABB}\SpecialCharTok{\textbackslash{}n}\StringTok{Timer, INTH}\SpecialCharTok{\textbackslash{}n}\StringTok{Keypad, TPU", fillcolor="\#B3E5FC"]}

\StringTok{  l1\_loop [label="Layer1 main loop}\SpecialCharTok{\textbackslash{}n}\StringTok{TDMA frame IRQ handler}\SpecialCharTok{\textbackslash{}n}\StringTok{L1CTL via sercomm", fillcolor="\#CE93D8"]}

\StringTok{  osmocon [label="osmocon connects}\SpecialCharTok{\textbackslash{}n}\StringTok{via pty → sercomm}\SpecialCharTok{\textbackslash{}n}\StringTok{→ /tmp/osmocom\_l2", fillcolor="\#FFE0B2"]}

\StringTok{  start {-}\textgreater{} load {-}\textgreater{} bss {-}\textgreater{} stacks {-}\textgreater{} ctors {-}\textgreater{} main {-}\textgreater{} hw\_init {-}\textgreater{} l1\_loop}
\StringTok{  l1\_loop {-}\textgreater{} osmocon [label="L1CTL}\SpecialCharTok{\textbackslash{}n}\StringTok{ready", style=dashed]}
\StringTok{\}\textquotesingle{}}
\FunctionTok{render\_graphviz}\NormalTok{(dot\_code)}
\end{Highlighting}
\end{Shaded}

\begin{center}\includegraphics[width=1.05in]{../../../../tmp/RtmpZyT17y/file2730a6c0b9cb9} \end{center}

\textbf{Point critique :} Le layer1 est chargé \textbf{directement} par
QEMU (\texttt{-kernel}), sans passer par le loader + osmocon upload. Le
protocole \texttt{osmoload\ memwrite} n'est pas implémenté dans
l'émulation UART, donc le chargement via loader échoue (le code reste
celui du loader à 0x820044).

\textbf{Diagnostic :} \texttt{xp\ /1xw\ 0x00820044} via le monitor QEMU
permet de distinguer les deux :

\begin{itemize}
\tightlist
\item
  \texttt{0xeb002876} → \textbf{layer1} (bl 0x82a224)
\item
  \texttt{0xeb000ea3} → \textbf{loader} (bl 0x823ad8)
\end{itemize}

\begin{center}\rule{0.5\linewidth}{0.5pt}\end{center}

\subsection{Orchestration --- tmux \&
Docker}\label{orchestration-tmux-docker}

\paragraph{Architecture de
déploiement}\label{architecture-de-duxe9ploiement}

\begin{Shaded}
\begin{Highlighting}[]
\NormalTok{dot\_code }\OtherTok{\textless{}{-}} \StringTok{\textquotesingle{}digraph deployment \{}
\StringTok{  graph [rankdir=TB, bgcolor=transparent, fontname=Helvetica, compound=true]}
\StringTok{  node [shape=box, style="filled,rounded", fontname=Helvetica, fontsize=10]}
\StringTok{  edge [fontname=Helvetica, fontsize=9]}

\StringTok{  subgraph cluster\_docker \{}
\StringTok{    label="Docker Container (openbsc{-}nitb)}\SpecialCharTok{\textbackslash{}n}\StringTok{{-}{-}privileged {-}{-}network=host"}
\StringTok{    style=filled; fillcolor="\#E8EAF6"; color="\#3F51B5"}

\StringTok{    subgraph cluster\_tmux \{}
\StringTok{      label="tmux session: gsm"; style=dashed; color="\#666"}

\StringTok{      w0 [label="Window 0: nitb}\SpecialCharTok{\textbackslash{}n}\StringTok{cd /etc/osmocom}\SpecialCharTok{\textbackslash{}n}\StringTok{osmo{-}nitb {-}{-}yes{-}i{-}really{-}}\SpecialCharTok{\textbackslash{}n}\StringTok{want{-}to{-}run{-}prehistoric{-}software",}
\StringTok{          fillcolor="\#BBDEFB"]}
\StringTok{      w1 [label="Window 1: calypso}\SpecialCharTok{\textbackslash{}n}\StringTok{cd /src/qemu}\SpecialCharTok{\textbackslash{}n}\StringTok{bash launch\_calypso.sh",}
\StringTok{          fillcolor="\#CE93D8"]}
\StringTok{      w2 [label="Window 2: trx}\SpecialCharTok{\textbackslash{}n}\StringTok{cd .../transceiver}\SpecialCharTok{\textbackslash{}n}\StringTok{./transceiver",}
\StringTok{          fillcolor="\#FFE0B2"]}
\StringTok{      w3 [label="Window 3: bts}\SpecialCharTok{\textbackslash{}n}\StringTok{cd /etc/osmocom}\SpecialCharTok{\textbackslash{}n}\StringTok{osmo{-}bts{-}trx {-}i 0.0.0.0",}
\StringTok{          fillcolor="\#C8E6C9"]}
\StringTok{    \}}

\StringTok{    entrypoint [label="entrypoint.sh}\SpecialCharTok{\textbackslash{}n}\StringTok{TUN setup (apn0)}\SpecialCharTok{\textbackslash{}n}\StringTok{ip\_forward}\SpecialCharTok{\textbackslash{}n}\StringTok{→ exec run.sh",}
\StringTok{                fillcolor="\#FFF9C4"]}

\StringTok{    tun [label="apn0}\SpecialCharTok{\textbackslash{}n}\StringTok{192.168.100.1/24}\SpecialCharTok{\textbackslash{}n}\StringTok{(GPRS NAT)", shape=ellipse, fillcolor="\#DCEDC8"]}
\StringTok{  \}}

\StringTok{  subgraph cluster\_host \{}
\StringTok{    label="Host Linux"; style=dashed; color="\#999"}
\StringTok{    wireshark [label="Wireshark}\SpecialCharTok{\textbackslash{}n}\StringTok{{-}i lo {-}k}\SpecialCharTok{\textbackslash{}n}\StringTok{{-}f udp port 4729}\SpecialCharTok{\textbackslash{}n}\StringTok{{-}Y gsmtap", fillcolor="\#E8F5E9"]}
\StringTok{    start\_sh [label="start.sh}\SpecialCharTok{\textbackslash{}n}\StringTok{wireshark \&}\SpecialCharTok{\textbackslash{}n}\StringTok{docker run ...", fillcolor="\#FFF9C4"]}
\StringTok{  \}}

\StringTok{  start\_sh {-}\textgreater{} wireshark [label="lance"]}
\StringTok{  start\_sh {-}\textgreater{} entrypoint [label="docker run"]}
\StringTok{  entrypoint {-}\textgreater{} tun [label="TUN/NAT"]}
\StringTok{  entrypoint {-}\textgreater{} w0 [label="run.sh", lhead=cluster\_tmux]}

\StringTok{  w0 {-}\textgreater{} w3 [style=invis]}
\StringTok{  w1 {-}\textgreater{} w2 [style=invis]}
\StringTok{\}\textquotesingle{}}
\FunctionTok{render\_graphviz}\NormalTok{(dot\_code)}
\end{Highlighting}
\end{Shaded}

\begin{center}\includegraphics[width=4.54in]{../../../../tmp/RtmpZyT17y/file2730ac004e8d} \end{center}

\paragraph{Séquence de démarrage
temporel}\label{suxe9quence-de-duxe9marrage-temporel}

\begin{Shaded}
\begin{Highlighting}[]
\NormalTok{dot\_code }\OtherTok{\textless{}{-}} \StringTok{\textquotesingle{}digraph timeline \{}
\StringTok{  graph [rankdir=TB, bgcolor=transparent, fontname=Helvetica]}
\StringTok{  node [shape=record, style=filled, fontname=Helvetica, fontsize=10]}
\StringTok{  edge [fontname=Helvetica, fontsize=9]}

\StringTok{  t0 [label="\{t=0s|start.sh|Wireshark → background}\SpecialCharTok{\textbackslash{}n}\StringTok{docker run → entrypoint.sh\}", fillcolor="\#FFF9C4"]}
\StringTok{  t1 [label="\{t=1s|entrypoint.sh|TUN apn0 + ip\_forward}\SpecialCharTok{\textbackslash{}n}\StringTok{exec run.sh\}", fillcolor="\#FFE0B2"]}
\StringTok{  t2 [label="\{t=2s|run.sh — Window 0|osmo{-}nitb démarre}\SpecialCharTok{\textbackslash{}n}\StringTok{Écoute Abis, GSUP\}", fillcolor="\#BBDEFB"]}
\StringTok{  t3 [label="\{t=5s|run.sh — Window 1|QEMU Calypso + osmocon}\SpecialCharTok{\textbackslash{}n}\StringTok{layer1 boot → L1CTL socket\}", fillcolor="\#CE93D8"]}
\StringTok{  t4 [label="\{t=8s|run.sh — Window 2|transceiver}\SpecialCharTok{\textbackslash{}n}\StringTok{Attend /tmp/osmocom\_l2 (30s max)}\SpecialCharTok{\textbackslash{}n}\StringTok{Puis bridge L1→TRX\}", fillcolor="\#FFE0B2"]}
\StringTok{  t5 [label="\{t=11s|run.sh — Window 3|osmo{-}bts{-}trx}\SpecialCharTok{\textbackslash{}n}\StringTok{Connecte au transceiver + nitb\}", fillcolor="\#C8E6C9"]}
\StringTok{  t6 [label="\{t=\textasciitilde{}15s|Stack opérationnelle|BTS radiating}\SpecialCharTok{\textbackslash{}n}\StringTok{GSMTAP → Wireshark\}", fillcolor="\#E8F5E9"]}

\StringTok{  t0 {-}\textgreater{} t1 {-}\textgreater{} t2 {-}\textgreater{} t3 {-}\textgreater{} t4 {-}\textgreater{} t5 {-}\textgreater{} t6}
\StringTok{\}\textquotesingle{}}
\FunctionTok{render\_graphviz}\NormalTok{(dot\_code)}
\end{Highlighting}
\end{Shaded}

\begin{center}\includegraphics[width=2.44in]{../../../../tmp/RtmpZyT17y/file2730a55659312} \end{center}

\begin{center}\rule{0.5\linewidth}{0.5pt}\end{center}

\subsection{Scripts de lancement}\label{scripts-de-lancement}

\paragraph{launch\_calypso.sh --- QEMU +
osmocon}\label{launch_calypso.sh-qemu-osmocon}

\begin{Shaded}
\begin{Highlighting}[]
\CommentTok{\#!/bin/bash}
\VariableTok{QEMU}\OperatorTok{=}\NormalTok{./build/qemu{-}system{-}arm}
\VariableTok{FIRMWARE}\OperatorTok{=}\VariableTok{$\{1}\OperatorTok{:{-}}\NormalTok{\textasciitilde{}/compal\_e88/layer1.highram.elf}\VariableTok{\}}
\VariableTok{MONSOCK}\OperatorTok{=}\NormalTok{/tmp/qemu{-}calypso{-}mon.sock}
\VariableTok{LOGFILE}\OperatorTok{=}\NormalTok{/tmp/qemu{-}calypso.log}
\FunctionTok{rm} \AttributeTok{{-}f} \StringTok{"}\VariableTok{$MONSOCK}\StringTok{"} \StringTok{"}\VariableTok{$LOGFILE}\StringTok{"}
\FunctionTok{killall} \AttributeTok{{-}q}\NormalTok{ qemu{-}system{-}arm }\DecValTok{2}\OperatorTok{\textgreater{}}\NormalTok{/dev/null}\KeywordTok{;} \FunctionTok{sleep}\NormalTok{ 0.3}

\VariableTok{$QEMU} \AttributeTok{{-}M}\NormalTok{ calypso }\DataTypeTok{\textbackslash{}}
  \AttributeTok{{-}kernel} \StringTok{"}\VariableTok{$FIRMWARE}\StringTok{"} \DataTypeTok{\textbackslash{}}
  \AttributeTok{{-}serial}\NormalTok{ pty }\DataTypeTok{\textbackslash{}}
  \AttributeTok{{-}monitor}\NormalTok{ unix:}\VariableTok{$MONSOCK}\NormalTok{,server,nowait }\DataTypeTok{\textbackslash{}}
  \OperatorTok{\textgreater{}}\StringTok{"}\VariableTok{$LOGFILE}\StringTok{"} \DecValTok{2}\OperatorTok{\textgreater{}\&}\DecValTok{1} \KeywordTok{\&}
\VariableTok{QEMU\_PID}\OperatorTok{=}\VariableTok{$!}

\CommentTok{\# Attendre le pty}
\VariableTok{PTS}\OperatorTok{=}\StringTok{""}
\ControlFlowTok{for}\NormalTok{ i }\KeywordTok{in} \VariableTok{$(}\FunctionTok{seq}\NormalTok{ 1 50}\VariableTok{)}\KeywordTok{;} \ControlFlowTok{do}
    \BuiltInTok{[} \OtherTok{{-}f} \StringTok{"}\VariableTok{$LOGFILE}\StringTok{"} \BuiltInTok{]} \KeywordTok{\&\&} \VariableTok{PTS}\OperatorTok{=}\VariableTok{$(}\FunctionTok{grep} \AttributeTok{{-}o} \StringTok{\textquotesingle{}/dev/pts/[0{-}9]*\textquotesingle{}} \StringTok{"}\VariableTok{$LOGFILE}\StringTok{"} \KeywordTok{|} \FunctionTok{head} \AttributeTok{{-}1}\VariableTok{)}
    \BuiltInTok{[} \OtherTok{{-}n} \StringTok{"}\VariableTok{$PTS}\StringTok{"} \BuiltInTok{]} \KeywordTok{\&\&} \ControlFlowTok{break}
    \FunctionTok{sleep}\NormalTok{ 0.1}
\ControlFlowTok{done}

\CommentTok{\# osmocon comme bridge L1CTL (pas d\textquotesingle{}upload, direct boot)}
\ExtensionTok{osmocon} \AttributeTok{{-}p} \StringTok{"}\VariableTok{$PTS}\StringTok{"} \AttributeTok{{-}m}\NormalTok{ c123xor }\KeywordTok{\&}
\VariableTok{OSMO\_PID}\OperatorTok{=}\VariableTok{$!}

\BuiltInTok{trap} \StringTok{"kill }\VariableTok{$QEMU\_PID}\StringTok{ }\VariableTok{$OSMO\_PID}\StringTok{ 2\textgreater{}/dev/null; rm {-}f }\VariableTok{$MONSOCK}\StringTok{; exit"}\NormalTok{ INT TERM}
\BuiltInTok{wait} \VariableTok{$QEMU\_PID}
\end{Highlighting}
\end{Shaded}

\paragraph{Commandes de diagnostic}\label{commandes-de-diagnostic}

\begin{Shaded}
\begin{Highlighting}[]
\CommentTok{\# Vérifier quel firmware est en RAM (loader vs layer1)}
\BuiltInTok{echo} \StringTok{"xp /1xw 0x00820044"} \KeywordTok{|} \ExtensionTok{socat} \AttributeTok{{-}}\NormalTok{ UNIX{-}CONNECT:/tmp/qemu{-}calypso{-}mon.sock}
\CommentTok{\# 0xeb002876 → layer1 | 0xeb000ea3 → loader}

\CommentTok{\# Dump des premières instructions}
\BuiltInTok{echo} \StringTok{"xp /16xw 0x00820000"} \KeywordTok{|} \ExtensionTok{socat} \AttributeTok{{-}}\NormalTok{ UNIX{-}CONNECT:/tmp/qemu{-}calypso{-}mon.sock}

\CommentTok{\# Vérifier que le CPU progresse}
\BuiltInTok{echo} \StringTok{"info registers"} \KeywordTok{|} \ExtensionTok{socat} \AttributeTok{{-}}\NormalTok{ UNIX{-}CONNECT:/tmp/qemu{-}calypso{-}mon.sock}

\CommentTok{\# Power button pulse}
\BuiltInTok{echo} \StringTok{"write 0xFFFE4800 0xDEAD"} \KeywordTok{|} \ExtensionTok{socat} \AttributeTok{{-}}\NormalTok{ UNIX{-}CONNECT:/tmp/qemu{-}calypso{-}mon.sock}

\CommentTok{\# Double power press}
\BuiltInTok{echo} \StringTok{"write 0xFFFE4800 0xBEEF"} \KeywordTok{|} \ExtensionTok{socat} \AttributeTok{{-}}\NormalTok{ UNIX{-}CONNECT:/tmp/qemu{-}calypso{-}mon.sock}

\CommentTok{\# Hexdump série (vérifier que layer1 envoie des trames sercomm)}
\FunctionTok{hexdump} \AttributeTok{{-}C} \StringTok{"}\VariableTok{$PTS}\StringTok{"} \KeywordTok{|} \FunctionTok{head} \AttributeTok{{-}20}
\end{Highlighting}
\end{Shaded}

\begin{center}\rule{0.5\linewidth}{0.5pt}\end{center}

\subsection{Stubs et périphériques non
implémentés}\label{stubs-et-puxe9riphuxe9riques-non-impluxe9mentuxe9s}

\begin{Shaded}
\begin{Highlighting}[]
\NormalTok{dot\_code }\OtherTok{\textless{}{-}} \StringTok{\textquotesingle{}digraph stubs \{}
\StringTok{  graph [rankdir=TB, bgcolor=transparent, fontname=Helvetica]}
\StringTok{  node [shape=box, style="filled,rounded", fontname=Helvetica, fontsize=10]}

\StringTok{  subgraph cluster\_impl \{}
\StringTok{    label="Implémentés"; style=filled; fillcolor="\#E8F5E9"}
\StringTok{    n1 [label="INTH (IRQ ctrl)", fillcolor="\#C8E6C9"]}
\StringTok{    n2 [label="Timer 1 \& 2", fillcolor="\#C8E6C9"]}
\StringTok{    n3 [label="UART Modem + IrDA", fillcolor="\#C8E6C9"]}
\StringTok{    n4 [label="SPI / ABB TWL3025", fillcolor="\#C8E6C9"]}
\StringTok{    n5 [label="I2C", fillcolor="\#C8E6C9"]}
\StringTok{    n6 [label="Keypad (stateful)", fillcolor="\#C8E6C9"]}
\StringTok{    n7 [label="TRX Bridge", fillcolor="\#C8E6C9"]}
\StringTok{    n8 [label="IRAM + Alias", fillcolor="\#C8E6C9"]}
\StringTok{  \}}

\StringTok{  subgraph cluster\_stub \{}
\StringTok{    label="Stubs (retourne 0)"; style=filled; fillcolor="\#FFF3E0"}
\StringTok{    s1 [label="0xFFFE6800 — tmr6800 (8{-}bit)", fillcolor="\#FFE0B2"]}
\StringTok{    s2 [label="0xFFFE8000 — mmio\_80xx (8{-}bit)", fillcolor="\#FFE0B2"]}
\StringTok{    s3 [label="0xFFFEF000 — conf (16{-}bit)", fillcolor="\#FFE0B2"]}
\StringTok{    s4 [label="0xFFFF9800 — mmio\_98xx (16{-}bit)", fillcolor="\#FFE0B2"]}
\StringTok{    s5 [label="0xFFFFF000 — DPLL (16{-}bit)", fillcolor="\#FFE0B2"]}
\StringTok{    s6 [label="0xFFFFF900 — Rhea bridge (16{-}bit)", fillcolor="\#FFE0B2"]}
\StringTok{    s7 [label="0xFFFFFB00 — CLKM (16{-}bit)", fillcolor="\#FFE0B2"]}
\StringTok{    s8 [label="0xFFFFFC00 — mmio\_fcxx (16{-}bit)", fillcolor="\#FFE0B2"]}
\StringTok{    s9 [label="0xFFFFFD00 — cntl (16{-}bit)", fillcolor="\#FFE0B2"]}
\StringTok{    s10 [label="0xFFFFFF00 — DIO (8{-}bit)", fillcolor="\#FFE0B2"]}
\StringTok{    s11 [label="0xFFF00000 — catch{-}all (1MB)", fillcolor="\#FFCCBC"]}
\StringTok{  \}}
\StringTok{\}\textquotesingle{}}
\FunctionTok{render\_graphviz}\NormalTok{(dot\_code)}
\end{Highlighting}
\end{Shaded}

\begin{center}\includegraphics[width=20.22in]{../../../../tmp/RtmpZyT17y/file2730a4a98a722} \end{center}

\begin{center}\rule{0.5\linewidth}{0.5pt}\end{center}

\subsection{Problèmes connus et
solutions}\label{probluxe8mes-connus-et-solutions}

\begin{longtable}[]{@{}
  >{\raggedright\arraybackslash}p{(\linewidth - 4\tabcolsep) * \real{0.3704}}
  >{\raggedright\arraybackslash}p{(\linewidth - 4\tabcolsep) * \real{0.2593}}
  >{\raggedright\arraybackslash}p{(\linewidth - 4\tabcolsep) * \real{0.3704}}@{}}
\toprule\noalign{}
\begin{minipage}[b]{\linewidth}\raggedright
Problème
\end{minipage} & \begin{minipage}[b]{\linewidth}\raggedright
Cause
\end{minipage} & \begin{minipage}[b]{\linewidth}\raggedright
Solution
\end{minipage} \\
\midrule\noalign{}
\endhead
\bottomrule\noalign{}
\endlastfoot
``Powering off due to keypress'' & Keypad initialisé à 0x00 (toutes
touches pressées) & Init \texttt{kp.keys\ =\ 0xFF} (active-low) \\
Layer1 jamais uploadé via osmocon & Protocole
\texttt{osmoload\ memwrite} non émulé dans UART & Charger layer1
directement via \texttt{-kernel} \\
\texttt{xp\ /16xi} → ``not supported on this arch'' & Désassemblage ARM
non supporté dans monitor & Utiliser \texttt{xp\ /16xw} (dump hex
brut) \\
Loader re-exécuté après \texttt{osmoload\ jump} & Code loader et layer1
à la même adresse 0x820000 & Comparer l'opcode à 0x820044 pour
distinguer \\
Flash init fail & Pas d'émulation flash NOR & Non bloquant, le firmware
continue \\
\end{longtable}

\begin{center}\rule{0.5\linewidth}{0.5pt}\end{center}

\subsection{Prochaines étapes}\label{prochaines-uxe9tapes}

\begin{Shaded}
\begin{Highlighting}[]
\NormalTok{dot\_code }\OtherTok{\textless{}{-}} \StringTok{\textquotesingle{}digraph roadmap \{}
\StringTok{  graph [rankdir=LR, bgcolor=transparent, fontname=Helvetica]}
\StringTok{  node [shape=box, style="filled,rounded", fontname=Helvetica, fontsize=10]}
\StringTok{  edge [fontname=Helvetica, fontsize=9]}

\StringTok{  done [label="Fait", shape=plaintext, fontsize=12, fontcolor="\#4CAF50"]}
\StringTok{  next [label="Prochaines}\SpecialCharTok{\textbackslash{}n}\StringTok{étapes", shape=plaintext, fontsize=12, fontcolor="\#FF9800"]}
\StringTok{  future [label="Long terme", shape=plaintext, fontsize=12, fontcolor="\#9C27B0"]}

\StringTok{  d1 [label="SoC complet}\SpecialCharTok{\textbackslash{}n}\StringTok{INTH+UART+SPI}\SpecialCharTok{\textbackslash{}n}\StringTok{+Timer+I2C", fillcolor="\#C8E6C9"]}
\StringTok{  d2 [label="TRX bridge}\SpecialCharTok{\textbackslash{}n}\StringTok{UDP 4729", fillcolor="\#C8E6C9"]}
\StringTok{  d3 [label="Keypad}\SpecialCharTok{\textbackslash{}n}\StringTok{stateful", fillcolor="\#C8E6C9"]}
\StringTok{  d4 [label="Boot direct}\SpecialCharTok{\textbackslash{}n}\StringTok{layer1 {-}kernel", fillcolor="\#C8E6C9"]}

\StringTok{  n1 [label="IRQ keypad}\SpecialCharTok{\textbackslash{}n}\StringTok{press/release", fillcolor="\#FFE0B2"]}
\StringTok{  n2 [label="Watchdog}\SpecialCharTok{\textbackslash{}n}\StringTok{timer", fillcolor="\#FFE0B2"]}
\StringTok{  n3 [label="osmoload}\SpecialCharTok{\textbackslash{}n}\StringTok{memwrite}\SpecialCharTok{\textbackslash{}n}\StringTok{(UART proto)", fillcolor="\#FFE0B2"]}
\StringTok{  n4 [label="GPRS}\SpecialCharTok{\textbackslash{}n}\StringTok{PCU+SGSN+GGSN", fillcolor="\#FFE0B2"]}

\StringTok{  f1 [label="Multi{-}MS}\SpecialCharTok{\textbackslash{}n}\StringTok{(N phones)", fillcolor="\#E1BEE7"]}
\StringTok{  f2 [label="Flash NOR}\SpecialCharTok{\textbackslash{}n}\StringTok{emulation", fillcolor="\#E1BEE7"]}
\StringTok{  f3 [label="Audio}\SpecialCharTok{\textbackslash{}n}\StringTok{osmo{-}gapk}\SpecialCharTok{\textbackslash{}n}\StringTok{ALSA", fillcolor="\#E1BEE7"]}
\StringTok{  f4 [label="Upstream}\SpecialCharTok{\textbackslash{}n}\StringTok{QEMU merge", fillcolor="\#E1BEE7"]}

\StringTok{  done {-}\textgreater{} d1 {-}\textgreater{} d2 {-}\textgreater{} d3 {-}\textgreater{} d4 [style=invis]}
\StringTok{  next {-}\textgreater{} n1 {-}\textgreater{} n2 {-}\textgreater{} n3 {-}\textgreater{} n4 [style=invis]}
\StringTok{  future {-}\textgreater{} f1 {-}\textgreater{} f2 {-}\textgreater{} f3 {-}\textgreater{} f4 [style=invis]}

\StringTok{  d4 {-}\textgreater{} n1 [style=dashed, color="\#999"]}
\StringTok{  n4 {-}\textgreater{} f1 [style=dashed, color="\#999"]}
\StringTok{\}\textquotesingle{}}
\FunctionTok{render\_graphviz}\NormalTok{(dot\_code)}
\end{Highlighting}
\end{Shaded}

\begin{center}\includegraphics[width=9.91in]{../../../../tmp/RtmpZyT17y/file2730a5bff088} \end{center}

\begin{center}\rule{0.5\linewidth}{0.5pt}\end{center}

\subsection{Annexe --- Fichiers source
QEMU}\label{annexe-fichiers-source-qemu}

\begin{longtable}[]{@{}
  >{\raggedright\arraybackslash}p{(\linewidth - 2\tabcolsep) * \real{0.6000}}
  >{\raggedright\arraybackslash}p{(\linewidth - 2\tabcolsep) * \real{0.4000}}@{}}
\toprule\noalign{}
\begin{minipage}[b]{\linewidth}\raggedright
Fichier
\end{minipage} & \begin{minipage}[b]{\linewidth}\raggedright
Rôle
\end{minipage} \\
\midrule\noalign{}
\endhead
\bottomrule\noalign{}
\endlastfoot
\texttt{hw/arm/calypso/calypso\_soc.c} & SoC principal, init de tous les
périphériques \\
\texttt{hw/arm/calypso/calypso\_uart.c} & UART 16550-like avec
chardev \\
\texttt{hw/arm/calypso/calypso\_inth.c} & Contrôleur d'interruptions 32
lignes \\
\texttt{hw/arm/calypso/calypso\_timer.c} & Timer avec auto-reload \\
\texttt{hw/arm/calypso/calypso\_spi.c} & SPI + registres ABB TWL3025 \\
\texttt{hw/arm/calypso/calypso\_i2c.c} & I2C stub \\
\texttt{hw/arm/calypso/calypso\_trx.c} & TRX bridge UDP (DSP API +
TPU) \\
\texttt{hw/arm/calypso/calypso\_machine.c} & Définition machine QEMU
\texttt{-M\ calypso} \\
\texttt{include/hw/arm/calypso/calypso\_soc.h} & Header SoC
(CalypsoSoCState) \\
\texttt{include/hw/arm/calypso/calypso\_trx.h} & Header TRX (IRQ
numbers, API) \\
\texttt{include/hw/arm/calypso/calypso\_uart.h} & Header UART
(TYPE\_CALYPSO\_UART) \\
\end{longtable}

\end{document}
